\documentclass[aspectratio=169]{beamer}

% ================== THEME & PACKAGES ==================
\usetheme[numbering=fraction, progressbar=foot]{metropolis}
\usepackage[utf8]{inputenc}
\usepackage[T1]{fontenc}
\usepackage[french]{babel}
\usepackage{lmodern}
\usepackage{amsmath, amssymb, mathtools}
\usepackage{booktabs}
\usepackage{hyperref}
\usepackage{xcolor}
\usepackage{tikz}
\usetikzlibrary{arrows.meta, positioning, calc, trees, shapes.geometric}
\usepackage{graphicx}
\usepackage{pgfplots}
\usepackage{tcolorbox}
\pgfplotsset{compat=1.17}

% ================== COLORS ==================
\definecolor{Primary}{RGB}{26,82,118}
\definecolor{Secondary}{RGB}{0,150,136}
\definecolor{Accent}{RGB}{45, 62, 80}
\definecolor{LightGray}{RGB}{245,245,245}

\setbeamercolor{frametitle}{fg=white, bg=Secondary}
\setbeamercolor{title}{fg=Primary}
\setbeamercolor{progress bar}{fg=Secondary, bg=Secondary!25}
\setbeamercolor{alerted text}{fg=Accent}
\setbeamercolor{structure}{fg=Primary}
\setbeamercolor{block title}{bg=Primary!90, fg=white}
\setbeamercolor{block body}{bg=Primary!10}

% ================== BOXES ==================
\newtcolorbox{keybox}[1]{
  colback=Secondary!5,
  colframe=Secondary!80,
  fonttitle=\bfseries\small,
  title=#1,
  boxrule=1.2pt,
  arc=2pt,
  left=3pt,
  right=3pt,
  top=3pt,
  bottom=3pt
}
\newtcolorbox{formulabox}{
  colback=LightGray,
  colframe=Primary,
  boxrule=0.8pt,
  arc=2pt,
  left=4pt,
  right=4pt,
  top=4pt,
  bottom=4pt
}

% ================== FONTS ==================
\setbeamerfont{frametitle}{size=\large}
\setbeamerfont{block title}{size=\small}

% ================== TITLE ==================
\title{\textbf{Modélisation et Valorisation d'Options}}
\subtitle{Black--Scholes • Binomial (CRR) • Trinomial (Boyle) • Américaines}
\author{Hassan EL QADI}
\institute{ENSA Agadir --- Finance \& Ingénierie de la Décision\\Encadré par M. Brahim El Asri}
\date{Octobre 2025}

% ================== SHORTCUTS ==================
\newcommand{\E}{\mathbb{E}}
\newcommand{\Q}{\mathbb{Q}}

\begin{document}

% ===== SLIDE 1: TITLE =====
\begin{frame}[plain]
  \titlepage
  \vspace{-8mm}
  \begin{center}\scriptsize
    Démo interactive : \href{https://options-binomial.vercel.app/pricing}{\textcolor{Secondary}{\textbf{options-binomial.vercel.app}}}
  \end{center}
\end{frame}

% ===== SLIDE 2: PLAN =====
\begin{frame}{Plan}
  \setbeamertemplate{section in toc}[sections numbered]
  \tableofcontents
\end{frame}

% ================== SECTION 1: INTRODUCTION ==================
\section{Introduction}

% ===== SLIDE 3: CONTEXTE, PROBLÉMATIQUE, OBJECTIFS =====
\begin{frame}{Contexte, Problématique et Objectifs}
  \small
  
  \textbf{Contexte}
  \begin{itemize}\scriptsize
    \item Les marchés dérivés jouent un rôle clé dans la gestion du risque, la spéculation et la couverture
  \end{itemize}

  \textbf{Problématique}
  \begin{itemize}\scriptsize
    \item Comment obtenir des \emph{prix justes} et des \emph{décisions d'exercice} robustes pour les options européennes et américaines ?
    \item Nécessité de méthodes rapides, stables et précises
  \end{itemize}

  \textbf{Objectifs}
  \begin{itemize}\scriptsize
    \item Implémenter BSM (benchmark), CRR (binomial) et Boyle (trinomial)
    \item Traiter l'\textbf{exercice anticipé} et calculer la frontière $S^*(t)$
    \item Valider via convergence, parité put-call et conditions de non-arbitrage
  \end{itemize}
\end{frame}

% ================== SECTION 2: RAPPELS ==================
\section{Rappels théoriques}

% ===== SLIDE 4: RAPPELS (1/2) =====
\begin{frame}{Définitions essentielles}
  \small
  
   \begin{columns}[T]
     \begin{column}{0.4\textwidth}
  \textbf{Paramètres de marché}
  \begin{itemize}\scriptsize
    \item $S_0$ : prix spot de l'actif sous-jacent
    \item $K$ : strike (prix d'exercice)
    \item $T$ : maturité (en années)
    \item $r$ : taux sans risque
    \item $q$ : taux de dividende continu
    \item $\sigma$ : volatilité
  \end{itemize}
  \end{column}
    \begin{column}{0.50\textwidth}
    \textbf{Payoffs}
  \begin{itemize}\scriptsize
    \item Call : $(S_T - K)^+$
    \item Put : $(K - S_T)^+$
  \end{itemize}
  \end{column}

  \end{columns}


  \vspace{2mm}
  \begin{keybox}{Mesure risque-neutre $\Q$}
    \scriptsize
    Sous la mesure risque-neutre : 
    $\E_\Q\!\left[\dfrac{S_{t+\Delta t}}{S_t}\right] = e^{(r-q)\Delta t}$
  \end{keybox}
\end{frame}

% ===== SLIDE 5: RAPPELS (2/2) =====
\begin{frame}{Parité et options américaines}
  \small
  
  \begin{formulabox}\scriptsize
    \textbf{Parité put-call européenne :}
    \vspace{1mm}
    
    $C - P = S_0e^{-qT} - Ke^{-rT}$
  \end{formulabox}
  
  \vspace{3mm}
  \begin{keybox}{Options américaines}
    \scriptsize
    \textbf{Principe :} À chaque nœud de l'arbre, on compare :
    \begin{itemize}\scriptsize
      \item Valeur intrinsèque (exercice immédiat)
      \item Valeur de continuation (attendre)
    \end{itemize}
    
    $V = \max(\text{intrinsèque}, \text{continuation})$
    
    \vspace{2mm}
    \textbf{Comportement :}
    \begin{itemize}\scriptsize
      \item \emph{Put américain} : exercice optimal quand $S$ est suffisamment bas
      \item \emph{Call américain} : exercice anticipé seulement avec dividendes importants
    \end{itemize}
  \end{keybox}
\end{frame}

% ================== SECTION 3: MODÈLES ==================
\section{Modèles de valorisation}

% ===== SLIDE 6: BLACK-SCHOLES-MERTON =====
\begin{frame}{Black--Scholes--Merton (Benchmark)}
  \small
  
  \begin{formulabox}\scriptsize
    \textbf{Prix d'un call européen :}
    
    $C = S_0e^{-qT}N(d_1) - Ke^{-rT}N(d_2)$
    
    \vspace{2mm}
    \textbf{Prix d'un put européen :}
    
    $P = Ke^{-rT}N(-d_2) - S_0e^{-qT}N(-d_1)$
    
    \vspace{2mm}
    où :
    
    $d_1 = \dfrac{\ln(S_0/K) + (r - q + \frac{1}{2}\sigma^2)T}{\sigma\sqrt{T}}$
    
    $d_2 = d_1 - \sigma\sqrt{T}$
  \end{formulabox}
  
  \vspace{3mm}
  \textbf{Pourquoi l'utiliser ?}
  \begin{itemize}\scriptsize
    \item Solution analytique fermée (calcul instantané)
    \item Calcul exact des grecques (Delta, Gamma, Vega, etc.)
    \item Sert de \textbf{référence} pour valider les méthodes numériques
    \item Largement utilisé par les praticiens
  \end{itemize}
\end{frame}

% ===== SLIDE 7: CRR BINOMIAL =====
% ===== SLIDE 7A: BINOMIAL PARAMETERS =====
\begin{frame}{Modèle binomial CRR — Paramètres}
  \small
  \textbf{Paramètres du modèle}
  \begin{formulabox}\scriptsize
    $\Delta t = \dfrac{T}{N}$ \quad (pas de temps)
    
    $u = e^{\sigma\sqrt{\Delta t}}$ \quad (facteur de hausse)
    
    $d = \dfrac{1}{u} = e^{-\sigma\sqrt{\Delta t}}$ \quad (facteur de baisse)
    
    $p = \dfrac{e^{(r-q)\Delta t} - d}{u - d}$ \quad (probabilité risque-neutre)
    
    $\text{disc} = e^{-r\Delta t}$ \quad (facteur d’actualisation)
    
    \vspace{1mm}
    Prix du sous-jacent au nœud $(i,j)$ :
    \[
      S_{i,j} = S_0 \cdot u^j \cdot d^{i-j}
    \]
  \end{formulabox}
\end{frame}



% ===== SLIDE 7B: BINOMIAL ALGORITHM =====
\begin{frame}{Modèle binomial CRR — Algorithme}
  \small
  \textbf{Étapes de calcul (Backward Induction)}
  \begin{enumerate}\scriptsize
    \item \textbf{Initialisation (maturité)} :\\
    Pour tous les nœuds $j$ à $i = N$, calculer :
    \[
      V_{N,j} = \text{payoff}(S_{N,j})
    \]
    Exemple : pour un call, $V_{N,j} = \max(S_{N,j} - K, 0)$.
    
    \vspace{2mm}
    \item \textbf{Rétropropagation (valeurs précédentes)} :\\
    Pour chaque $i = N-1, N-2, \ldots, 0$ et chaque $j \in [0, i]$ :
    \[
      C_{i,j} = \text{disc} \cdot [p \cdot V_{i+1,j+1} + (1-p) \cdot V_{i+1,j}]
    \]
    
    \vspace{2mm}
    \item \textbf{Condition selon le type d’option :}
    \begin{itemize}
      \item \textbf{Européenne :} $V_{i,j} = C_{i,j}$
      \item \textbf{Américaine :} $V_{i,j} = \max(\text{valeur intrinsèque}(S_{i,j}), C_{i,j})$
    \end{itemize}
    
    \vspace{2mm}
    \item \textbf{Prix initial :} $V_0 = V_{0,0}$
  \end{enumerate}
\end{frame}

\begin{frame}{Modèle binomial CRR — Schéma de l'algorithme}
  \small
  \begin{center}
    \begin{tikzpicture}[
      scale=0.75,
      node distance=2cm and 3cm,
      price/.style={circle, draw=black!80, thick, minimum size=0.8cm, font=\tiny},
      value/.style={rectangle, draw=green!70, fill=green!10, thick, minimum size=0.7cm, font=\tiny},
      payoff/.style={rectangle, draw=red!70, fill=red!10, thick, minimum size=0.7cm, font=\tiny},
      arrow/.style={->, >=stealth, thick},
      label/.style={font=\tiny, text=black}
    ]
    
    % Timeline
    \draw[->, thick] (-0.6, -3.8) -- (6, -3.8);
    \node[font=\tiny] at (0, -4) {$t=0$};
    \node[font=\tiny] at (1.8, -4) {$t=\Delta t$};
    \node[font=\tiny] at (4.2, -4) {$t=2\Delta t = T$};
    
    % Nodes at t=0
    \node[price] (S0) at (0,0) {$S_0$};
    \node[value, below=0.3cm of S0] (V0) {$V_{0,0}$};
    
    % Nodes at t=1
    \node[price] (S11) at (2,1.3) {$S_{1,1}$};
    \node[label, above=0.05cm of S11] {$S_0 u$};
    \node[value, right=0.25cm of S11] (V11) {$V_{1,1}$};
    
    \node[price] (S10) at (2,-1.3) {$S_{1,0}$};
    \node[label, below=0.05cm of S10] {$S_0 d$};
    \node[value, right=0.25cm of S10] (V10) {$V_{1,0}$};
    
    % Nodes at t=2 (maturity)
    \node[price] (S22) at (4.2,2.4) {$S_{2,2}$};
    \node[label, above=0.05cm of S22] {$S_0 u^2$};
    \node[payoff, right=0.25cm of S22] (P22) {$V_{2,2}$};
    
    \node[price] (S21) at (4.2,0) {$S_{2,1}$};
    \node[label, right=0.05cm of S21, xshift=1cm] {$S_0 ud$};
    \node[payoff, right=0.25cm of S21] (P21) {$V_{2,1}$};
    
    \node[price] (S20) at (4.2,-2.4) {$S_{2,0}$};
    \node[label, below=0.05cm of S20] {$S_0 d^2$};
    \node[payoff, right=0.25cm of S20] (P20) {$V_{2,0}$};
    
    % Price tree connections (gray)
    \draw[gray, thick] (S0) -- (S11) node[midway, above left, font=\tiny] {$u$};
    \draw[gray, thick] (S0) -- (S10) node[midway, below left, font=\tiny] {$d$};
    \draw[gray, thick] (S11) -- (S22) node[midway, above left, font=\tiny] {$u$};
    \draw[gray, thick] (S11) -- (S21) node[midway, below left, font=\tiny] {$d$};
    \draw[gray, thick] (S10) -- (S21) node[midway, above left, font=\tiny] {$u$};
    \draw[gray, thick] (S10) -- (S20) node[midway, below left, font=\tiny] {$d$};
    
    % Backward induction arrows
    % From t=2 to t=1
    \draw[arrow, red!70, dashed] (P22) to[bend left=10] node[midway, above, font=\tiny] {$p$} (V11);
    \draw[arrow, red!70, dashed] (P21) to[bend right=8] node[midway, below, font=\tiny] {$1-p$} (V11);
    
    \draw[arrow, red!70, dashed] (P21) to[bend left=8] node[midway, above, font=\tiny] {$p$} (V10);
    \draw[arrow, red!70, dashed] (P20) to[bend right=10] node[midway, below, font=\tiny] {$1-p$} (V10);
    
    % From t=1 to t=0
    \draw[arrow, green!70, very thick] (V11) to[bend right=15] node[midway, above left, font=\tiny] {$p$} (V0);
    \draw[arrow, green!70, very thick] (V10) to[bend left=15] node[midway, below left, font=\tiny] {$1-p$} (V0);
    
    % Step labels (moved further right with more spacing)
    \node[draw, fill=red!10, text width=3cm, align=left, font=\tiny, thick] at (12, 2.2) {
      \textbf{Étape 1 :} \\[1mm]
      Payoffs à maturité $T$ \\[0.5mm]
      $V_{2,j} = \text{payoff}(S_{2,j})$ \\[0.5mm]
      Ex: Call $\max(S_{2,j}-K, 0)$
    };
    
    \node[draw, fill=orange!10, text width=4.5cm, align=left, font=\tiny, thick] at (12, 0) {
      \textbf{Étape 2 :} \\[1mm]
      Rétropropagation \\[0.5mm]
      $C_{i,j} = e^{-r\Delta t}
      \times [p \cdot V_{i+1,j+1}
      + (1-p) \cdot V_{i+1,j}]$
    };
    
    \node[draw, fill=green!10, text width=2.5cm, align=left, font=\tiny, thick] at (12, -2.2) {
      \textbf{Étape 3 :} \\[1mm]
      Prix initial de l'option \\[0.5mm]
      $V_0 = V_{0,0}$ \\[0.5mm]
      Résultat final
    };
    
    \end{tikzpicture}
  \end{center}
  
  \vspace{-3mm}
  
  \begin{formulabox}\scriptsize
    \textbf{Condition selon le type :} \quad
    Européenne : $V_{i,j} = C_{i,j}$ \quad | \quad
    Américaine : $V_{i,j} = \max(\text{payoff}(S_{i,j}), C_{i,j})$
  \end{formulabox}
\end{frame}


% ===== SLIDE 8A: TRINOMIAL PARAMETERS =====
\begin{frame}{Modèle trinomial de Boyle — Paramètres}
  \small
  \textbf{Paramètres du modèle}
  \begin{formulabox}\scriptsize
    $u = e^{\sigma\sqrt{2\Delta t}}$, \quad $m = 1$, \quad $d = \dfrac{1}{u}$

    $a = e^{(r-q)\Delta t/2}$, \quad $b = e^{\sigma\sqrt{\Delta t/2}}$

    \vspace{1mm}
    \textbf{Probabilités risque-neutres :}
    \[
      p_u = \left(\dfrac{a - b^{-1}}{b - b^{-1}}\right)^2, \quad
      p_d = \left(\dfrac{b - a}{b - b^{-1}}\right)^2, \quad
      p_m = 1 - p_u - p_d
    \]
    
    \vspace{1mm}
    Indexation centrée : \quad
    $S_{i,k} = S_0 \cdot u^k$ \quad avec $k \in [-i, \ldots, i]$
    
    \vspace{1mm}
    Facteur d’actualisation : \quad $\text{disc} = e^{-r\Delta t}$
  \end{formulabox}
\end{frame}

% ===== SLIDE 8B: TRINOMIAL ALGORITHM =====
\begin{frame}{Modèle trinomial de Boyle — Algorithme}
  \small
  \textbf{Étapes de calcul (Backward Induction)}
  \begin{enumerate}\scriptsize
    \item \textbf{Initialisation à maturité :}\\
    Pour tous les $k \in [-N, N]$ :
    \[
      V_{N,k} = \text{payoff}(S_{N,k})
    \]
    Exemple : $V_{N,k} = \max(S_{N,k} - K, 0)$ pour un call.
    
    \vspace{2mm}
    \item \textbf{Rétropropagation :}\\
    Pour $i = N-1, N-2, \ldots, 0$ et $k \in [-i, i]$ :
    \[
      C_{i,k} = \text{disc} \cdot [p_u V_{i+1,k+1} + p_m V_{i+1,k} + p_d V_{i+1,k-1}]
    \]
    
    \vspace{2mm}
    \item \textbf{Condition selon le type d’option :}
    \begin{itemize}
      \item \textbf{Européenne :} $V_{i,k} = C_{i,k}$
      \item \textbf{Américaine :} $V_{i,k} = \max(\text{intrinsèque}(S_{i,k}), C_{i,k})$
    \end{itemize}
    
    \vspace{2mm}
    \item \textbf{Prix initial :} $V_0 = V_{0,0}$
  \end{enumerate}
\end{frame}

% ===== SLIDE: TRINOMIAL BOYLE SCHEMA =====
\begin{frame}{Modèle trinomial de Boyle — Schéma de l'algorithme}
  \small
  
  \begin{center}
    \begin{tikzpicture}[
      scale=0.65,
      node distance=1.8cm and 2.5cm,
      price/.style={circle, draw=black!80, thick, minimum size=0.65cm, font=\tiny},
      value/.style={rectangle, draw=green!70, fill=green!10, thick, minimum size=0.6cm, font=\tiny},
      payoff/.style={rectangle, draw=red!70, fill=red!10, thick, minimum size=0.6cm, font=\tiny},
      arrow/.style={->, >=stealth, thick},
      label/.style={font=\tiny, text=black}
    ]
    
    % Timeline
    \draw[->, thick] (-0.6, -4.5) -- (7, -4.5);
    \node[font=\tiny] at (0, -4.75) {$t=0$};
    \node[font=\tiny] at (3, -4.75) {$t=\Delta t$};
    \node[font=\tiny] at (6.2, -4.75) {$t=2\Delta t = T$};
    
    % Nodes at t=0
    \node[price] (S0) at (0,0) {$S_0$};
    \node[value, below=0.25cm of S0] (V0) {$V_{0,0}$};
    
    % Nodes at t=1
    \node[price] (S12) at (3,2) {$S_{1,2}$};
    \node[label, above=0.05cm of S12] {$S_0 u$};
    \node[value, right=0.4cm of S12] (V12) {$V_{1,2}$};
    
    \node[price] (S11) at (3,0) {$S_{1,1}$};
    \node[label, right=0.05cm of S11, xshift=1.4cm] {$S_0$};
    \node[value, right=0.4cm of S11] (V11) {$V_{1,1}$};
    
    \node[price] (S10) at (3,-2) {$S_{1,0}$};
    \node[label, below=0.05cm of S10] {$S_0 d$};
    \node[value, right=0.4cm of S10] (V10) {$V_{1,0}$};
    
    % Nodes at t=2 (maturity)
    \node[price] (S24) at (6.2,3.2) {$S_{2,4}$};
    \node[label, above=0.05cm of S24] {$S_0 u^2$};
    \node[payoff, right=0.8cm of S24] (P24) {$V_{2,4}$};
    
    \node[price] (S23) at (6.2,1.6) {$S_{2,3}$};
    \node[label, right=0.05cm of S23, xshift=1.4cm] {};
    \node[payoff, right=0.8cm of S23] (P23) {$V_{2,3}$};
    
    \node[price] (S22) at (6.2,0) {$S_{2,2}$};
    \node[label, right=0.05cm of S22, xshift=1.4cm] {};
    \node[payoff, right=0.8cm of S22] (P22) {$V_{2,2}$};
    
    \node[price] (S21) at (6.2,-1.6) {$S_{2,1}$};
    \node[label, right=0.05cm of S21, xshift=1.4cm] {};
    \node[payoff, right=0.8cm of S21] (P21) {$V_{2,1}$};
    
    \node[price] (S20) at (6.2,-3.2) {$S_{2,0}$};
    \node[label, below=0.05cm of S20] {$S_0 d^2$};
    \node[payoff, right=0.8cm of S20] (P20) {$V_{2,0}$};
    
    % Price tree connections (gray) - from t=0 to t=1
    \draw[gray, thick] (S0) -- (S12) node[midway, above left, font=\tiny] {$u$};
    \draw[gray, thick] (S0) -- (S11) node[midway, above, font=\tiny, xshift=-2mm] {$m$};
    \draw[gray, thick] (S0) -- (S10) node[midway, below left, font=\tiny] {$d$};
    
    % Price tree connections (gray) - from t=1 to t=2
    \draw[gray, thick] (S12) -- (S24) node[midway, above left, font=\tiny] {$u$};
    \draw[gray, thick] (S12) -- (S23) node[midway, below left, font=\tiny, yshift=1mm] {$m$};
    \draw[gray, thick] (S12) -- (S22) node[midway, below left, font=\tiny] {$d$};
    
    \draw[gray, thick] (S11) -- (S23) node[midway, above left, font=\tiny] {$u$};
    \draw[gray, thick] (S11) -- (S22) node[midway, above, font=\tiny, yshift=1mm] {$m$};
    \draw[gray, thick] (S11) -- (S21) node[midway, below left, font=\tiny] {$d$};
    
    \draw[gray, thick] (S10) -- (S22) node[midway, above left, font=\tiny] {$u$};
    \draw[gray, thick] (S10) -- (S21) node[midway, below left, font=\tiny, yshift=-1mm] {$m$};
    \draw[gray, thick] (S10) -- (S20) node[midway, below left, font=\tiny] {$d$};
    
    % Backward induction arrows - from t=2 to t=1
    \draw[arrow, red!70, dashed] (P24) to[bend left=8] node[midway, above, font=\tiny,xshift=3mm] {$p_u$} (V12);
    \draw[arrow, red!70, dashed] (P23) -- (V12) node[midway, right, font=\tiny, xshift=3mm] {$p_m$};
    \draw[arrow, red!70, dashed] (P22) to[bend right=8] node[midway, below, font=\tiny, xshift=3mm] {$p_d$} (V12);
    
    \draw[arrow, red!70, dashed] (P23) to[bend left=6] node[midway, above, font=\tiny, xshift=3mm] {$p_u$} (V11);
    \draw[arrow, red!70, dashed] (P22) -- (V11) node[midway, right, font=\tiny, xshift=3mm] {$p_m$};
    \draw[arrow, red!70, dashed] (P21) to[bend right=6] node[midway, below, font=\tiny, xshift=3mm] {$p_d$} (V11);
    
    \draw[arrow, red!70, dashed] (P22) to[bend left=8] node[midway, above, font=\tiny, xshift=3mm] {$p_u$} (V10);
    \draw[arrow, red!70, dashed] (P21) -- (V10) node[midway, right, font=\tiny, xshift=3mm] {$p_m$};
    \draw[arrow, red!70, dashed] (P20) to[bend right=8] node[midway, below, font=\tiny, xshift=3mm] {$p_d$} (V10);
    
    % From t=1 to t=0
    \draw[arrow, green!70, very thick] (V12) to[bend right=10] node[midway, above left, font=\tiny] {$p_u$} (V0);
    \draw[arrow, green!70, very thick] (V11) -- (V0) node[midway, left, font=\tiny, xshift=-1mm] {$p_m$};
    \draw[arrow, green!70, very thick] (V10) to[bend left=10] node[midway, below left, font=\tiny] {$p_d$} (V0);
    
    % Step labels
    \node[draw, fill=red!10, text width=3cm, align=left, font=\tiny, thick] at (15, 2.5) {
      \textbf{Étape 1 :} \\[1mm]
      Payoffs à maturité $T$ \\[0.5mm]
      $V_{2,j} = \text{payoff}(S_{2,j})$ \\[0.5mm]
      Ex: Call $\max(S_{2,j}-K, 0)$
    };
    
    \node[draw, fill=orange!10, text width=3cm, align=left, font=\tiny, thick] at (15, 0) {
      \textbf{Étape 2 :} \\[1mm]
      Rétropropagation \\[0.5mm]
      $C_{i,j} = e^{-r\Delta t} \times$ \\[0.5mm]
      $[p_u \cdot V_{i+1,j+2} + p_m \cdot V_{i+1,j+1}$ \\[0.5mm]
      $+ p_d \cdot V_{i+1,j}]$
    };
    
    \node[draw, fill=green!10, text width=3cm, align=left, font=\tiny, thick] at (15, -2.5) {
      \textbf{Étape 3 :} \\[1mm]
      Prix initial de l'option \\[0.5mm]
      $V_0 = V_{0,0}$ \\[0.5mm]
      Résultat final
    };
    
    \end{tikzpicture}
  \end{center}
  
  \vspace{-3mm}
  
  \begin{formulabox}\scriptsize
    \textbf{Condition :} Européenne : $V_{i,j} = C_{i,j}$ \quad | \quad
    Américaine : $V_{i,j} = \max(\text{payoff}(S_{i,j}), C_{i,j})$
  \end{formulabox}
\end{frame}

% ================== SECTION 4: OPTIONS AMÉRICAINES ==================
\section{Options américaines}

% ===== SLIDE 9: FRONTIÈRE D'EXERCICE =====
\begin{frame}{Frontière d'exercice anticipé}
  \small
  
  \begin{keybox}{Définition}
    \scriptsize
    La \textbf{frontière d'exercice} $S^*(t)$ sépare deux zones :
    \begin{itemize}\scriptsize
      \item \textbf{Zone d'exercice} : pour un put, si $S \le S^*(t)$ $\Rightarrow$ exercer immédiatement
      \item \textbf{Zone de continuation} : si $S > S^*(t)$ $\Rightarrow$ attendre est optimal
    \end{itemize}
  \end{keybox}
  \begin{center}
  \begin{tikzpicture}
    \begin{axis}[
      width=0.75\textwidth,
      height=0.35\textwidth,
      xlabel={\scriptsize Temps $t$},
      ylabel={\scriptsize Prix critique $S^*(t)$},
      xmin=0, xmax=1,
      ymin=60, ymax=110,
      grid=both,
      tick label style={font=\scriptsize},
      legend style={font=\scriptsize, fill=white},
      legend pos=north west
    ]
      \addplot[Secondary, thick, mark=*] coordinates{
        (0.2,62) (0.4,66) (0.6,72) (0.8,82) (0.95,100)
      };
      \addlegendentry{Frontière put américain}
      \addplot[dashed, Primary] coordinates{(0,90) (1,90)};
      \addlegendentry{Strike $K=90$}
    \end{axis}
  \end{tikzpicture}
  \end{center}
  
  {\scriptsize\textit{Note :} Au début, la frontière peut être absente si l'exercice n'est jamais optimal (continuation domine).}
\end{frame}

% ===== SLIDE 10: ALGORITHME FRONTIÈRE =====
\begin{frame}{Calcul de la frontière d'exercice --- Algorithme}
  \small
  
  \textbf{Étapes de l'algorithme}
  \begin{enumerate}\scriptsize
    \item Construire l'arbre complet des prix du sous-jacent
    \item Initialiser les valeurs terminales avec les payoffs
    \item \textbf{Remonter l'arbre} (backward induction) :
    \begin{itemize}\scriptsize
      \item Calculer $C$ (valeur de continuation)
      \item Calculer $V = \max(\text{intrinsèque}, C)$ pour option américaine
    \end{itemize}
    \item \textbf{À chaque date $t_i$} : identifier les nœuds où l'exercice est optimal
    \begin{itemize}\scriptsize
      \item Put : garder le \emph{plus grand} $S$ exercé $\Rightarrow$ $S^*(t_i)$
      \item Call : garder le \emph{plus petit} $S$ exercé $\Rightarrow$ $S^*(t_i)$
    \end{itemize}
    \item Stocker les points $(t_i, S^*(t_i))$ et tracer la courbe $S^*(t)$
  \end{enumerate}
\end{frame}

% ================== SECTION 5: RÉSULTATS ==================
\section{Résultats et validation}

% ===== SLIDE 11: CONVERGENCE =====
\begin{frame}{Convergence --- Précision vs complexité}
  \small
  \begin{itemize}\scriptsize
    \item \textbf{CRR (binomial)} : Erreur $\sim O(1/\sqrt{N})$
    \item \textbf{Boyle (trinomial)} : Erreur $\sim O(1/N)$ à $N$ comparable
    \item \textbf{Complexité} : $O(N^2)$ en temps, $O(N)$ en mémoire
  \end{itemize}

  \begin{center}
    \begin{tikzpicture}
      \begin{axis}[
        width=0.75\textwidth,
        height=0.39\textwidth,
        xlabel={\scriptsize Nombre de pas $N$},
        ylabel={\scriptsize Erreur relative (\%)},
        xmin=0, xmax=120,
        ymin=0, ymax=1.0,
        grid=both,
        legend pos=north east,
        tick label style={font=\scriptsize},
        legend style={font=\scriptsize, fill=white}
      ]
        \addplot[thick, mark=*, Primary] coordinates{
          (10,0.90) (20,0.55) (40,0.32) (60,0.25) (80,0.21) (100,0.18)
        };
        \addlegendentry{CRR (binomial)}
        
        \addplot[thick, mark=square*, Secondary] coordinates{
          (10,0.62) (20,0.34) (40,0.19) (60,0.15) (80,0.13) (100,0.12)
        };
        \addlegendentry{Boyle (trinomial)}
      \end{axis}
    \end{tikzpicture}
  \end{center}
\end{frame}

% ===== SLIDE 12: TESTS DE VALIDATION =====
\begin{frame}{Tests de validation}
  \small
  
  \begin{columns}[T]
    \begin{column}{0.50\textwidth}
      \begin{keybox}{Arbitrage et parité}
        \scriptsize
        \textbf{Parité put-call européenne :}
        
        $C - P = S_0e^{-qT} - Ke^{-rT}$
        
        Erreur observée : $< 0.02\%$
        
        \vspace{2mm}
        \textbf{Bornes d'arbitrage :}
        
        $C \ge \max(S_0e^{-qT} - Ke^{-rT}, 0)$
        
        $P \ge \max(Ke^{-rT} - S_0e^{-qT}, 0)$
      \end{keybox}
      
      \vspace{2mm}
      \begin{keybox}{Mesure risque-neutre}
        \scriptsize
        \textbf{Probabilités valides :}
        
        Binomial : $p \in [0,1]$
        
        Trinomial : $p_u, p_m, p_d \in [0,1]$
        
        \vspace{1mm}
        \textbf{Martingale :}
        
        $\E_\Q[S_{t+\Delta t}] = S_t e^{(r-q)\Delta t}$
      \end{keybox}
    \end{column}
    
    \begin{column}{0.48\textwidth}
    \vspace{1.4cm}
      \begin{block}{Américaines}\scriptsize
        \textbf{Prime d'exercice anticipé :}
        
        Put américain : Prime $\ge 0$
        
        Call américain ($q=0$) : Prime $= 0$
        
        \vspace{2mm}
        \textbf{Frontière $S^*(t)$ :}
        
        Monotonicité quasi croissante (avec tolérance pour discrétisation)
      \end{block}
    \end{column}
  \end{columns}
\end{frame}

% ================== SECTION 6: ARCHITECTURE ==================
\section{Architecture et démonstration}

% ===== SLIDE 13: ARCHITECTURE LOGICIELLE =====
\begin{frame}{Architecture logicielle}
  \small
  
  \vspace{5mm}
  
  \textbf{Backend (Python)}
  \begin{itemize}\scriptsize
    \item Framework : FastAPI (API REST haute performance)
    \item Librairies : NumPy, SciPy (calculs numériques)
    \item Fonctionnalités : pricing, calcul des grecques, frontières d'exercice, tests de validation
  \end{itemize}
  
  \vspace{-1.5mm}
  \textbf{Frontend (TypeScript)}
  \begin{itemize}\scriptsize
    \item Framework : Next.js (React)
    \item Styling : Tailwind CSS
    \item Visualisation : Recharts (graphiques interactifs)
  \end{itemize}

  
  \vspace{-1.5mm}
  \begin{keybox}{Flux de données}
    \scriptsize
    Interface utilisateur $\leftrightarrow$ API REST $\leftrightarrow$ Moteur de pricing (BS/CRR/Boyle)
  \end{keybox}
\end{frame}

% ===== SLIDE 14: DÉMONSTRATION =====
\begin{frame}{Démonstration de la plateforme}
  \small
  \begin{center}
      \vspace{15mm}
        {\huge \href{https://options-binomial.vercel.app}{\textcolor{Secondary}{\textbf{options-binomial.vercel.app}}}}
  \end{center}
\end{frame}

% ================== SECTION 7: CONCLUSION ==================
\section{Conclusion}

% ===== SLIDE 15: CONCLUSION ET PERSPECTIVES =====
\begin{frame}{Conclusion et perspectives}
  \small
  
  \vspace{2mm}
  \textbf{Réalisations}
  \begin{itemize}\scriptsize
    \item Implémentation complète de BSM, CRR et Boyle
    \item Traitement des options américaines avec calcul de la frontière d'exercice
    \item Validation rigoureuse : convergence, parité put-call, conditions de non-arbitrage
    \item Application web interactive connectée aux données de marché réelles
    \item Documentation complète et tests unitaires
  \end{itemize}
  
  \vspace{-2mm}
  \textbf{Extensions futures}
  \begin{itemize}\scriptsize
    \item \textbf{Options bermudanes} : exercice à dates spécifiques
    \item \textbf{Surface de volatilité} : calibration et modélisation de la volatilité implicite
    \item \textbf{Options exotiques} : barrières, asiatiques, lookback
  \end{itemize}
\end{frame}

% ===== SLIDE 16: MERCI =====
\begin{frame}[plain]
  \begin{center}
    \vspace{15mm}
    {\Huge \textbf{Merci pour votre attention !}}
    
    \vspace{10mm}
    {\Large \textbf{Questions ?}}
    
    \vspace{12mm}
    {\small \textcolor{Primary}{\textbf{hassanelqadi3@gmail.com}}}
    
    \vspace{2mm}
    {\small Démo interactive : \href{https://options-binomial.vercel.app}{\textcolor{Secondary}{\textbf{options-binomial.vercel.app}}}}
  \end{center}
\end{frame}

\end{document}